\cotefile{64}
\scanpage{66}
4. Détermination de $M_{11}(4)$.

Soit $E$ une courbe elliptique sur $S$ (à car. résiduelles $\ne 2$), considérons
\[
M={}_4E \tag*{(38)}
\]
schéma en groupes sur $S$ loc isom. à $(\mathbb Z/4\mathbb Z)^2$, i.e. \deleted{syst. local} faisceau abélien \deleted{\unsure{qui|pour}} \inserted{sur} $S_{\text{ét}}$ qui est un $(\mathbb Z/4\mathbb Z)_S$ module loc. libre de rg 2.

\deleted{(39)} de sorte qu'on a
\diagram{\begin{tikzcd}[ampersand replacement=\&,column sep=small,row sep=small]
\text{(39)} \& 0 \arrow[r] \& M_0 \arrow[r] \& M \arrow[r,"\mathrm{can}"] \& M_0 \arrow[r] \& 0 \\
\& \& M \arrow[u,"\mathrm{can}"'] \arrow[ur,"2\,\mathrm{id}_M"']
\end{tikzcd}}
où
\[
M_0=M\otimes_{\mathbb Z/4\mathbb Z}\mathbb F_2\simeq {}_2E .
\]

La donnée d'un $M$ équivaut à celle d'un $M_0$, (i.e. d'un $\underline{J}$ rev. étale de degré 3 de $S$ \deleted{\unsure{($\underline J$}}
\[
\underline J\simeq M_0^\ast,\qquad
M_0=V_{\underline J}(\mathbb F_2)\simeq
\mathbb F_2^{\underline J}/\underbrace{\mathbb F_2}_{\mathrm{diag}}\big), \tag*{(40)}
\]
\uline{et} d'un torseur \inserted{$P$} sous
\[
\unsure{N|V}_{\underline J}(M_0)\simeq\underline{\mathrm{End}}_{\mathbb F_2,S}(M_0)
\simeq M_0\oplus\mathbb F_2^{\underline{\unsure{\omega|w}}} \tag*{(41)}
\]
% slanted
\margin{$\underline\omega$ est le rev. étale de degré 2 ``des orientations'' associé au rev. $\underline J$ de degré 3 de $S$}

Un épinglage de $M={}_4E$ (i.e. \deleted{d'un} la donnée d'un iso
\[
(\mathbb Z/4\mathbb Z)^2_S\simeq M\;(={}_4E) \tag*{(42)}
\]
équivaut à celle d'une \inserted{trivialisation de $\underline J$, plus d'une} section \unsure{(trivialisation)|(après translation)}
% archive sheet number: 65

\scanpage{67}\mspage{12}
\deleted{plus celle d'une} de ce torseur, \deleted{qui} \inserted{Ce torseur} peut aussi s'interpréter \deleted{comme le torseur des sections} \inserted{ainsi : la structure d'extension (39)} définit la structure d'extension
\diagram{\begin{tikzcd}[ampersand replacement=\&,column sep=small,row sep=0.3em]
\& \& \underline{\mathrm{End}}(M_0) \& \& \underline{\mathrm{End}}(M_0) \\
\& \& \rotatebox{90}{$\simeq$} \& \& \rotatebox{90}{$\simeq$} \\
\text{(43)} \& \unsure{0'|0} \arrow[r] \& V_{\underline{J}}(M_0) \arrow[r] \& V_{\underline J}(M) \arrow[r] \& V_J(M_0) \arrow[r] \& 1
\end{tikzcd}}
où pour \unsure{tt|le|$A$} schéma \inserted{(ou faisceau)} en groupes comm. $A$ sur $S$, on pose
\[
V_{\underline J}(A)\simeq \mathrm{Ker}\big(A^{\underline J}\xrightarrow{\mathrm{somme}} A\big) \tag*{(44)}
\]
(faisceau loc isom. à $A^2$), et $V_{\underline J}$ est un foncteur exact en $A$, commutant bien sûr au chgt. de bas. Pour $A$ annulé par 2, on a un isom
\[
V_{\underline J}(A)\simeq\mathrm{Hom}(M_0,A)\qquad\text{si }2\,\mathrm{id}_A=0 \tag*{(45)}
\]
dû au fait que l'on a
\[
M_0\simeq\mathbb F_2^{\underline J}/\underbrace{\mathbb F_2}_{\mathrm{diag}} \tag*{(46)}
\]

Donc à une section \deleted{\unsure{de|d}} canonique partout non nulle de $V_J(M_0)$, correspondant à l'application d'inclusion $\underline J\simeq M_0^\ast\hookrightarrow M_0$, ou encore à la section \deleted{\unsure{``iden|``idem|évidente}} $\mathrm{id}_{M_0}$ de $\underline{\mathrm{End}}(M_0)$, dont l'image inverse dans (43) est le torseur cherché.
% archive sheet number: 66

\scanpage{68}
\deleted{Un épinglage de $M$ s'identifie donc à une}

On peut considérer, indépendamment de la question d'épingler $J$, \deleted{celle de trouver une section \unsure{?|de|s}} \inserted{l'épinglage ``relatif'' \unsure{consistant|consistera}} en \deleted{une} la donnée d'une section du torseur $P_M$ sous $\underline{\mathrm{End}}(M_0)$, ce qui équivaut \inserted{\unsure{du|donc}} aussi à la donnée d'un morphisme
\[
\underline J\longrightarrow\deleted{\unsure{V_{\underline J}(}}M\deleted{\unsure{)|\}}} \tag*{(47)}
\]
qui relève le morphisme d'inclusion
\[
\underline J\hookrightarrow M_0
\]
et qui soit une section de $V_{\underline J}(M)$ i.e. qui soit à trace nulle (rel. à $S$).

\unsure{On a|Ceci est|Ceci} en fait automatiquement un morphisme
\[
\underline J\longrightarrow M\setminus M_0
\]
donc en composant avec \deleted{l'inclusion}
\[
(\underline J\dashrightarrow)\,M\hookrightarrow E\longrightarrow\Sigma_{\underline J}\simeq E/(1,\sigma_E)
\]
on trouve un morphisme
\[
\underline J\longrightarrow\Delta_t \tag*{(48)}
\]
et la condition de relèvement pour (47) signifie aussi que le composé de (48)
% archive sheet number: 67

\scanpage{69}\mspage{13}
avec le morphisme canonique (étale de degré 2)
\[
f_t:\ \Delta_t\longrightarrow\underline{J}
\]
soit l'identité -- i.e. (48) définit une \uline{face} du fibré en octaèdres combinatoires $\Delta_t$ sur $S$. Du coup l'octaèdre est $\underline J$-épinglé, en d'autres termes, il est can. isomorphe au \deleted{$\unsure{\underline J\text{-}}$}\unsure{octàèdre|octaèdre} standart défini par $\underline J$, \inserted{et une section de $\mu_4^\ast$,} en prenant le revêtement biquadratique $\Sigma_{\underline J,2}$ de $\Sigma_{\underline J}$ \ldots
% archive sheet number: 68

\scanpage{70}
III

Épinglage de Legendre des courbes elliptiques
% archive sheet number: 69

\scanpage{71}
\uline{Digression} Épinglages de revêtements par \deleted{\unsure{cadres|cadre} fibrés \unsure{tangts|tangents}} rev. Kummériens des \deleted{\unsure{f|fi}}espaces tangents (cas relatif)
\par\noindent\rule{\linewidth}{0.4pt}\par

Soient
\[
 f:X\longrightarrow S \tag*{(1)}
\]
un morphisme lisse de dim relative 1, $t$ une section de $X$ sur $S$, $X^\ast=X\setminus t(S)$.

On considère des schémas \inserted{$X'$} finis sur $\deleted{S} X$, lisses sur $S$ de dim relative 1, munis d'une section $t'$ au dessus de $t$, et ``modérément ramifiés'' en $t'$ sur $X$, avec un degré de ramification $n$ -- ce qui signifie que
% slanted
\margin{\uline{NB} $n$ est premier aux car. résiduelles de $S$ !!!}
loc pour la top. étale sur $X,X'$, $X'$ est isom au voisinage de tt point $x'$ de $t'(S)$ à un $X$-schéma Kummérien de $X$ au voisinage de $x=p(x')$, de la forme
\margin{$p:X'\longrightarrow X$}
\[
 X'_0=\mathrm{Spec}\,\underline O_X[z]/(z^n-\varphi)
\]
où $\varphi$ est une équation locale du diviseur relatif $t(S)\subset X$ au voisinage de $x$.

Sous ces conditions, via l'hom norme locale, on définit un isom. canonique
\[
 T_{X'/\deleted{t}S,t'}^{\otimes n}\xrightarrow{\overset{k_{t'}}{\sim}}T_{X/S,t} \tag*{(2)}
\]
% archive sheet number: 70

\scanpage{72}
Si d'autre part $\underline G$ est un schéma en groupes fini étale sur $S$, et si $X'$ est muni d'une structure de $\underline G_X$-torseur au dessus de $X^\ast$, alors le stabilisateur
% slanted
\margin{\uline{NB} On suppose que $\underline G$ opère sur $X'$, de façon que $X'^\ast=X'|X^\ast$ devienne un $\underline G_X$-torseur sur $X^\ast$}
dans $\underline G_{\deleted{\unsure{X|S}}}$ \deleted{\unsure{du}} de la section $t'$ est un sous-schéma en groupes \inserted{$\underline G_{t'}$} de $\underline G$, \deleted{\unsure{comme groupe|comme grou}} \inserted{qui} \deleted{\unsure{image dans|ment à|munit de}} \inserted{opère \inserted{fid.} sur} $T_{X'/S,t'}$ \uline{et} dont l'image dans $\deleted{T_{X'/S,t'}}\,\mathrm{Gm}_S\simeq\underline{\mathrm{Aut}}_{O_S}T_{X'/S,t'}$ est $\mu_{n\,S}$, d'où un monomorphisme canonique
\[
\mu_{n\,S}\hookrightarrow\underline G . \tag*{(3)}
\]

Par ailleurs, désignant par
\[
f^\circ:X^\ast\longrightarrow S
\]
le morphisme induit par $f$, le $\underline G_{X^\ast}$-torseur $X'^\ast$ définit une section
\[
\xi_{X'/S}\in H^0(S,R^1 f^\circ_\ast f^{\circ\ast}(\underline G)) . \tag*{(4)}
\]

Tout ceci posé, on a :

\uline{Proposition} \quad $(X',t')$ au dessus de $(X,\unsure{t|f})$, avec l'opération de $\underline G$ dessus, est défini à isom unique près par la connaissance de (3) (4) et de (2). \deleted{\unsure{(ces données|(ces derniers|(ce dernier}} De façon précise,
% archive sheet number: 71

\scanpage{73}
\deleted{\unsure{pour une section|pour ces sections}} considérons (pour $S,X,t$ fixés) la catégorie des $(X',t',\underline G)$, \deleted{où $X'$ est un $X$-schéma} \inserted{tels que dessus, avec} comme morphismes
\[
\unsure{u:|u<}\ (X',t',G)\longrightarrow(X'_1,t'_1,G_1)
\]
% slanted
\margin{\uline{NB} On n'exige pas que $n=n_1$. Pour $X',t',G$ fixés, $(X'_1,t'_1,G_1)$ est déterminé à isom can. près, quand on se donne $G\to G_1$}
les \deleted{\unsure{hom.|homs}} couples \inserted{de} $X'\to X'_1$ ($X$-morphismes et de $G\to G_1$, compatibles dans un sens évident avec les sections marquées et les opérations. Considérons d'autre part les \deleted{triples} systèmes
\[
(\underline G,\deleted{\unsure{\xi}}n,\deleted{\unsure{\xi}}i,\deleted{\unsure{n|i}}\xi,\deleted{\unsure{k}}T,k)
\]
où $n\in\mathbb N^\ast$, \inserted{$\underline G$ un schéma en groupes fini étale sur $S$,} $\xi\in H^0(S,R^1f^\circ_\ast f^{\circ\ast}(\underline G))$ tels que les rev. locaux (\unsure{ét. loc.|étale.|étales} sur $S$) de $X$ qu'elle décrit soient modérément ramifiés d'ordre $n$ le long de $t$, $i:\mu_{n\,S}\hookrightarrow\underline G$ est un mono. de groupes, $T$ un Module inv. sur $S$, et $k$ un isom
\[
T^{\otimes n}\xrightarrow{\sim} T_{X/S,t}
\]
avec la définition ad hoc des morphismes (\uline{NB} \unsure{bien|bête|bien entendu} un $\underline G\to\underline G_1$ doit induire un \unsure{épimorphisme|homomorphisme|\deleted{épi}morphisme} $\mu_{n\,S}\to\mu_{n_1\,S}$, qui soit l'hom \deleted{\unsure{gén.|groupe|qui}} $\lambda\longmapsto\lambda^{n/n_1}$ \ldots). On a un foncteur canonique de la première catégorie dans la
% archive sheet number: 72

\scanpage{74}
seconde, et je dis que ce foncteur est une équivalence de catégories.

Le cas qui m'intéresse est celui où $\underline G=\mu_2\unsure{\simeq|\xrightarrow{\sim}|\underset{\mathrm{can}}{\simeq}}\{\pm1\}_S$, et où $n=2$ -- donc où on s'intéresse aux rev. quadratiques $X'$ de $X/S$, étales au dessus de $X$ et ``totalement ramifiés le long de la section $t$'', avec $S$ à car. rés $\ne 2$. \deleted{\unsure{Si|Soit|Ici}} Pour une section $\xi\in H^0(S,R^1f_\circ^\ast(\mu_{2\,X^\ast}))$ fixée, satisfaisant à la condition de ramification totale, la catégorie des \deleted{\unsure{ce|ces|cou}}revêtements \inserted{correspondant à cette $\xi$} est donc équivalente à celle des \deleted{\unsure{rev.|rev}} couples $(T,k)$ d'un Module inv. $T$ sur $S$, et d'un isom
\[
T^{\otimes2}\simeq T_{X,t} .
\]

Le cas qui nous intéresse pour les courbes elliptiques est celui où $X$ est de la forme $\Sigma_{J}^\ast$ (fibré en droites projectives associé à un rev. étale cubique $J$ de $S$), \inserted{déduit d'une courbe elliptique $X'$ comme $X'/\{\pm1\}$,} et où $t$ est la section de $\Sigma_J^\ast/S$ \deleted{\unsure{quotient|qui}} image de la section nulle de $X'=E$ \unsure{en|un}
% archive sheet number: 73

\scanpage{75}\mspage{1}
Soit $E$ une courbe elliptique -- on travaille sur un corps alg. clos $k$ pour simplifier le propos, mais la seule hypothèse essentielle est que la base soit à car. résiduelles $\ne2$. On désigne par \deleted{\unsure{$Z$|$J$|$\Sigma$}}
\[
J={}_2E(k)^\ast \tag*{(1)}
\]
l'ens des pts de $E$ d'ordre 2, de sorte que l'on a
\[
E/(\pm\mathrm{id}_E)\simeq\Sigma_J\qquad
\text{\deleted{\unsure{quadrique de|quotient de|spectre de}} \inserted{droite projective} $J$-marquée.} \tag*{(2)}
\]
\[
\Sigma_J\simeq\mathbb P(V_J(k))\qquad
0\mathbin{-} V_J(k)\to k^J\xrightarrow{\mathrm{somme}} k\to 0 \tag*{(3)}
\]

La donnée de $E$ est équivalente à celle de $J$ i.e. de $\Sigma_J$ muni du diviseur étale $D_J$ de degré 3, plus \unsure{le|les} point
\[
t\in\Sigma_J^\ast=\Sigma_J\setminus D_J \tag*{(4)}
\]
image de l'élément $0_E$, plus une droite $T$ \inserted{sur $k$} (savoir $T_{E,0}$), munie d'un isomorphisme
\[
T^{\otimes2}\xrightarrow{\sim}T_{\Sigma,t} \tag*{(5)}
\]
(\deleted{\unsure{Réinterprétation|Réinterpréter|isomorphisme|isomorphismes}} \inserted{équivalence} de catégories-groupoïdes -- valable en fait sur base $S$ quelconque à car. rés. $\ne2$, et compatible avec changement de base \ldots).

1 \quad Notons qu'on a un isomorphisme canonique
\[
T_\Sigma\simeq\underline O_\Sigma(2)(\underline{\unsure{\omega|w}}) \tag*{(6)}
\]
où $\underline\omega=\underline\omega(J)$ est l'ens. des deux ordres circulaires sur $J$, et où $L(\underline\omega)$ (pour un objet tel que faisceau abélien etc) désigne $L\wedge_{\{\pm1\}}\underline\omega$.
% archive sheet number: 74

\scanpage{76}\mspage{2}
On a un candidat naturel (pour $\Sigma_{\deleted{\unsure{J|I}}}\inserted{=\Sigma_J}$, $\deleted{D_J}\inserted{D=D_J}$, $t\in\Sigma^\ast$ donnés) comme droite $T$ munie de (5), savoir
\[
T_0\simeq O_t(1) \tag*{\deleted{(7)}}
\]
à cela près qu'il y a \uline{deux} isom. naturels
\[
T_0^{\otimes2}=O_t(2)\simeq T_{\Sigma,t}\simeq O_t(\unsure{2|r})(\underline{\unsure{\omega|w}}),\tag*{\deleted{(8)}}
\]
correspondant aux deux éléments $w$ de $\underline\omega$, savoir
\[
\deleted{\unsure{t}}\ \xi\longmapsto\xi\wedge w \qquad
T_0^{\otimes2}\xrightarrow{\sim}O_t(2)\simeq T_{\Sigma,t}\simeq\underline O_t(2)(\underline\omega) \tag*{(7)}
\]

On va lever cette ambiguité, en remplaçant $T_0=O_t(1)$ par une variante tordue. Le tordre par un torseur sous $\mu_2$ \inserted{$=\mu_2(k)$} ne nous avance pas, car $\deleted{\unsure{O|T}}T_0^{\otimes2}$ n'aura pas changé (à isom can. près) on va le tordre par un torseur (qu'on précisera) $Q_0$ sous $\mu_4=\mu_4(k)$ (qui peut donc s'interpréter comme l'ens des sommets d'un \uline{carré}, dont $\mu_4$ est le groupe des rotations, avec $\{i,-i\}=\mu_4^\ast$ comme \uline{paire} de gén. privilégiés \ldots).

On posera donc
\[
T_0=\underline O_t(1)\wedge_{\mu_4}Q_0 \tag*{(8)}
\]
d'où
\[
T_0^{\otimes2}\simeq\underline O_t(2)\wedge_{\mu_2}D_{Q_0} \tag*{(9)}
\]
où
\[
D_{Q_0}=Q_0/\{\pm1\} \tag*{(10)}
\]
est l'ens des deux diagonales du carré $Q_0$. Supposons donc qu'on se donne un carré \inserted{combinatoire (``sommets'')} $Q_0$, avec des isom. \deleted{\unsure{Aut|Art} \unsure{$D_{Q_0}$|$Q_0$} \unsure{$\Sigma$|$Z$|$\simeq$} \unsure{entre|noter} le groupe des}
\[
\left\{\begin{aligned}
\underline\omega\,Q_0\quad &\text{(ens des orientations de $Q_0$)} &&\simeq\mu_4^\ast=\{i,-i\}\\
D_{Q_0}\quad &\text{(ens des diagonales de $Q_0$)} &&\simeq\underline\omega
\end{aligned}\right. \tag*{(11)}
\]
% archive sheet number: 75

\scanpage{77}
Alors, définissant $T_0$ par (8) (pour $t\in\Sigma_{J}^\ast$ donné) on a un isom canonique
\[
T_0^{\otimes2}\simeq T_{\Sigma,t}\tag*{(12)}
\]

\deleted{\unsure{une fois|ceci posé|en fait} \illegible{2} \unsure{supposé donné un} \unsure{$Q_0$} \unsure{$J$|de} \illegible{1} \unsure{qui soit|quand}} \inserted{Ceci posé, on appellera \uline{épinglage de Legendre}} de $E$, subordonné au carré $Q_0$ (muni des données (11)) \deleted{\unsure{tout}} \inserted{un} isomorphisme
\[
T=T_{E,0}\xleftarrow{\overset{k}{\sim}}T_0 \tag*{(13)}
\]
compatible avec les hom. norme, i.e. rendant commutatif le diagramme
\diagram{\begin{tikzcd}[ampersand replacement=\&,column sep=small,row sep=large]
\text{(14)} \& T_0^{\otimes2}\arrow[rr,"{\overset{k^{\otimes2}}{\sim}}"] \arrow[dr,"{(12)}"] \arrow[d,phantom,"{\unsure{!}}"] \& \& T^{\otimes2}\arrow[dl,"{(5)}"]\\
\& {} \& T_{\Sigma,t} \& {}
\end{tikzcd}}

Il y en exactement deux, différant par le signe i.e. formant un torseur sous $\mu_2$.

Bien entendu, on veut pouvoir définir la notion d'épinglage de $E$ (en termes d'un $Q_0$ donné une fois pour \unsure{toutes|toute}) pour $E$ variable -- mais alors $\underline{\unsure{\omega|w}}=\underline\omega_J$ n'est pas encore déterminé (puisque $J$ ne l'est pas). Donc la donnée d'un épinglage de Legendre sera défini comme le couple \inserted{$(\alpha,k)$} formé d'un isom.
\[
\alpha:\underline\omega_J\xleftarrow{\sim}D_{Q_0}=Q_0/\{\pm1\} \tag*{(15)}
\]
\uline{et} d'un isom. $k$ (13) satisfaisant la compatibilité (14),
% archive sheet number: 76
